% Capítulo incluído por tcc-ufv.tex com \include: sem preâmbulo e sem \begin{document}.
\chapter{Metodologia}\label{cap:metodologia}

Os requisitos estão no \autoref{qua:requisitos}, e o \autoref{lst:media}
mostra a função que calcula a média das avaliações.

% quadro e codigo são floats da classe, com listas próprias
\begin{quadro}[htbp]
\begin{ilustracao}[\linewidth]{\caption{Requisitos funcionais do protótipo}%
                               \label{qua:requisitos}}%
                              {\elaboracaopropria}
  \renewcommand{\arraystretch}{1.25}
  \begin{tabularx}{\linewidth}{|c|X|c|}
    \hline
    \textbf{ID} & \textbf{Descrição} & \textbf{Prioridade} \\
    \hline
    RF01 & Cadastrar usuário com nome, \emph{e-mail} e senha. & Alta \\
    \hline
    RF02 & Exportar as avaliações em \acs{csv}. & Média \\
    \hline
  \end{tabularx}
\end{ilustracao}
\end{quadro}

\begin{codigo}[htbp]
\begin{ilustracao}[\linewidth]{\caption{Média das avaliações}\label{lst:media}}%
                              {\elaboracaopropria}
\begin{lstlisting}[language=Python]
def media(notas):
    """Devolve a média; 0 para a lista vazia."""
    return sum(notas) / len(notas) if notas else 0.0
\end{lstlisting}
\end{ilustracao}
\end{codigo}
